\SourcePage{561}
\SGAsection{6}{附录：局部自由层的行列式计算}

\textbf{6.1。}先考虑任意\Term{ringed-topos}{赋环拓扑斯} $X=(C,O)$。对于\Term{constant-sheaf}{常值层} $\mathbf Z$ 的每个正\Term{section}{截面} $\underline n\in\Gamma(\mathbf Z_X)$，以及每个 $O$-\Term{module}{模} $M\in\operatorname{Mod}_O(X)$，下列 $O$-\Term{module}{模}的定义是清楚的：
\SourcePage{562}
\[
 M^{\underline n},\qquad \bigotimes^{\underline n}M,\qquad
 \bigwedge^{\underline n}M,\qquad\text{等等。}
 \tag{6.1.1}
\]
\Term{functor}{函子} $(\ )^{\underline n}$、$\bigotimes^{\underline n}$ 和 $\bigwedge^{\underline n}$ 将\Term{free-module}{自由模}（相应地，\Term{locally-free-module}{局部自由模}）变为\Term{free-module}{自由模}（相应地，\Term{locally-free-module}{局部自由模}）。

更\Term{generic}{一般}地，考虑 $\mathbf Z_X$ 的\Term{section}{截面} $\underline n_1,\ldots,\underline n_r$，以及一个有 $r$ 个变量的\Term{multifunctor}{多函子}（为明确起见，取\Term{covariant-functor}{协变}的情形）
\[
 T:\operatorname{Mod}^{\underline n_1}(O)\times\cdots\times
 \operatorname{Mod}^{\underline n_r}(O)\longrightarrow\operatorname{Mod}(O),
\]
其中 $\operatorname{Mod}^{\underline n}_O(X)$ 表示 $X$ 上由\Term{rank}{秩}为 $\underline n$ 的\Term{locally-free-module}{局部自由} $O$-\Term{module}{模}构成的\Term{fibred-category}{纤维范畴}。令
\[
 M=T(O^{\underline n_1},\ldots,O^{\underline n_r})
\]
以及
\[
 G=\operatorname{GL}(\underline n_1)\times\cdots\times
 \operatorname{GL}(\underline n_r),\qquad
 \operatorname{GL}(\underline n_i)=\operatorname{Aut}_O(O^{\underline n_i}).
\]
立即可知，$T$ 定义一个\Term{representation}{表示}
\[
 \alpha_T:G\longrightarrow\operatorname{Aut}_O(M).\tag{6.1.2}
\]
显然，若 $M$ 本身\Term{locally-free-module}{局部自由}且\Term{rank}{秩}为 $\underline m$，则对每个复合\Term{functor}{函子} $T'T$，其中
\[
 T':\operatorname{Mod}^{\underline m}(O)\longrightarrow\operatorname{Mod}(O),
\]
都有
\[
 \alpha_{T'T}=\alpha_{T'}\circ\alpha_T.\tag{6.1.3}
\]

反过来，给定\Term{representation}{表示} $\alpha=\alpha_T:G\to\operatorname{Aut}_O(M)$，对于任意一列\Term{rank}{秩}分别为 $\underline n_1,\ldots,\underline n_r$ 的\Term{locally-free-module}{局部自由} $O$-\Term{module}{模}
\SourcePage{563}
$E_1,\ldots,E_r$，可借助\Term{representation}{表示} $\alpha$，在\Term{local-scheme}{局部}\Term{trivialization}{平凡化}的坐标片上将\Term{free-module}{自由} $O$-\Term{module}{模} $T(O^{\underline n_1},\ldots,O^{\underline n_r})$ \Term{gluing}{粘合}，从而恢复 $T(E_1,\ldots,E_r)$。确切地说，序列 $E_1,\ldots,E_r$ 对应于一个右 $G$-\Term{torsor}{挠子} $P=P_{E_1,\ldots,E_r}$，并有\Term{canonical}{典范}\Term{isomorphism}{同构}
\[
 E_i\xrightarrow{\sim}P\mathbin{\wedge^G}O^{\underline n_i}
 \qquad(i=1,\ldots,r).\tag{6.1.3}
\]
这里 $\mathbin{\wedge^G}$ 表示\Term{contracted-product}{缩积}（produit contracté），$G$ 通过第 $i$ 个\Term{projection}{投影} $G\to\operatorname{GL}(\underline n_i)$ \Term{action}{作用}于 $O^{\underline n_i}$。于是，每个\Term{representation}{表示}
\[
 \alpha:G\longrightarrow\operatorname{Aut}_O(M)
\]
都使 $G$ 通过 $O$-\Term{automorphism}{自同构}\Term{action}{作用}于 $O^{\underline m}$，因而定义一个\Term{multifunctor}{多函子} $T=T_\alpha$，使得
\[
 T_\alpha(E_1,\ldots,E_r)=P_{E_1,\ldots,E_r}\mathbin{\wedge^G}M.
 \tag{6.1.4}
\]
此外，对于每个复合\Term{representation}{表示} $\alpha=\alpha_1\circ\alpha_2$，都有
\[
 T_\alpha=T_{\alpha_1}\circ T_{\alpha_2}.\tag{6.1.5}
\]

\textbf{6.2。}对于每个\Term{rank}{秩}为 $\underline n$ 的\Term{locally-free-module}{局部自由} $O$-\Term{module}{模} $E$，以 $\det(E)$ 表示\Term{invertible-module}{可逆} $O$-\Term{module}{模} $\bigwedge^{\underline n}E$。沿用 6.1 的记号，\Term{functor}{函子} $\det(T(E_1,\ldots,E_r))$ 由相应的\Term{representation}{表示}
\[
 \det\circ\alpha_T:G\longrightarrow\mathbf G_m=\operatorname{GL}(1)
 \tag{6.2.1}
\]
确定，精确到\Term{canonical}{典范}\Term{isomorphism}{同构}。

现在设 $X$ 是\Term{ringed-site}{赋环景} $(\operatorname{Sch}/S)$ 配备\Term{etale-topology}{平展拓扑}后，其上\Term{sheaf}{层}所成的\Term{ringed-topos}{赋环拓扑斯}。于是可利用下面的命题计算 $\det(T(E_1,\ldots,E_r))$，或\Term{equivalence}{等价}地，计算\Term{representation}{表示} $\det\circ\alpha_T$。

\SourcePage{564}
\begin{SGAresult}{命题}{6.2.2}
对于 $\mathbf Z$ 的每个正\Term{section}{截面} $\underline n$，$\operatorname{GL}(\underline n)$ 的\Term{commutator}{换位子}所成的\Term{etale-sheaf}{平展群层} $\operatorname{Der}(\operatorname{GL}(\underline n))$ 由
\[
 \operatorname{SL}(\underline n)=
 \operatorname{Ker}\bigl(\det:\operatorname{GL}(\underline n)\longrightarrow\mathbf G_m\bigr)
\]
\Term{representable-functor}{表示}。
\end{SGAresult}

\begin{SGAresult}{推论}{6.2.3}
每个\Term{representation}{表示}
\[
 \alpha:G=\operatorname{GL}(\underline n_1)\times\cdots\times
 \operatorname{GL}(\underline n_r)\longrightarrow\mathbf G_m\tag{6.2.4}
\]
都形如
\[
 \alpha(Y)(s)=\bigl(\det(s_1)\bigr)^{\underline m_1}\cdots
 \bigl(\det(s_r)\bigr)^{\underline m_r},\tag{6.2.5}
\]
其中 $Y\in\operatorname{Ob}(\operatorname{Sch}/S)$，$s=(s_1,\ldots,s_r)\in G(Y)$，而 $\underline m_1,\ldots,\underline m_r$ 是由 $\alpha$ 唯一确定的 $\mathbf Z$ 的\Term{section}{截面}。
\end{SGAresult}

\begin{SGAresult}{推论}{6.2.6}
沿用 6.1 的记号，有
\[
 \det\bigl(T_\alpha(E_1,\ldots,E_r)\bigr)=
 \det(E_1)^{\otimes\underline m_1}\otimes\cdots\otimes
 \det(E_r)^{\otimes\underline m_r},
\]
其中 $\underline m_1,\ldots,\underline m_r$ 是由 $\alpha$ 唯一确定的 $\mathbf Z$ 的\Term{section}{截面}。
\end{SGAresult}

\textbf{注记 6.2.7。}为确定 (6.2.5) 和 (6.2.6) 中的指数 $\underline m_1,\ldots,\underline m_r$，可以在 (6.2.5) 中依次代入
\[
 s_i=\lambda\cdot\operatorname{id}\quad(\lambda\in\mathbf G_m(Y)),
 \qquad s_j=\operatorname{id}_{\operatorname{GL}(\underline n_j)(Y)}
 \quad(j\ne i;\ i,j=1,\ldots,r),
\]
从而
\[
 \det(s_i)=\lambda^{\underline n_i},\qquad
 \alpha(Y)(s)=\lambda^{\underline n_i\underline m_i}.
\]
例如，这样可由 (6.2.6) \Term{pushout}{推出}
\[
 \det(E_1\otimes E_2)=\bigl(\det(E_1)\bigr)^{\otimes\underline n_2}
 \otimes\bigl(\det(E_2)\bigr)^{\otimes\underline n_1},\tag{6.2.8}
\]
\[
 \det(E_1\oplus E_2)=\det(E_1)\otimes\det(E_2),\tag{6.2.9}
\]
以及
\[
 \det\left(\bigwedge^{\underline m}E_1\right)=
 \bigl(\det(E_1)\bigr)^{\binom{\underline n_1-1}{\underline m-1}},
 \tag{6.2.10}
\]
等等。

\SourcePage{565}
\textbf{6.3。}先说明如何由 6.2.2 \Term{pushout}{推出} 6.2.3。显然，只须对
\[
 s_{(i)}=(1,\ldots,s_i,\ldots,1)\qquad(i=1,\ldots,r)
\]
证明 (6.2.5)，所以可设 $r=1$、$G=\operatorname{GL}(\underline n_1)$。由 6.2.2，$\alpha$ 通过 $\operatorname{GL}(\underline n_1)/\operatorname{SL}(\underline n_1)$ \Term{resolution}{分解}，即有\Term{commutative-diagram}{交换图}
\[
\begin{tikzcd}[column sep=large,row sep=large]
\operatorname{GL}(\underline n_1) \arrow[r,"\alpha"] \arrow[d,"j"'] & \mathbf G_m\\
\operatorname{GL}(\underline n_1)/\operatorname{SL}(\underline n_1) \arrow[ur,"\beta"'] &
\end{tikzcd}
\]
其中 $j$ 是\Term{canonical}{典范}\Term{epimorphism}{满态射}，$\beta$ 是从 $\operatorname{GL}(\underline n_1)/\operatorname{SL}(\underline n_1)$ 到 $\mathbf G_m$ 的一个\Term{representation}{表示}。但\Term{homomorphism}{同态}
\[
 \det:\operatorname{GL}(\underline n_1)\longrightarrow\mathbf G_m
\]
有一个\Term{section}{截面}，所以由 $\det$ 诱导的\Term{monomorphism}{单态射}
\[
 \operatorname{GL}(\underline n_1)/\operatorname{SL}(\underline n_1)
 \longrightarrow\mathbf G_m
\]
是\Term{isomorphism}{同构}。另一方面，已知（SGA 3，第 VIII 讲 1.3）$\mathbf G_m$ 到自身的每个\Term{representation}{表示}都形如
\[
 s\longmapsto s^{\underline m},
\]
其中 $\underline m$ 是 $\mathbf Z$ 的一个\Term{section}{截面}。断言由此得证。

推论 6.2.6 只是 6.2.3 的另一种表述。（注意，6.2.3 也可直接由 SGA 3，第 XXII 讲 6.2.1 \Term{pushout}{推出}；但这里希望通过 6.2.2 给出一个初等证明。）

\SourcePage{566}
\textbf{6.4。6.2.2 的证明。}只须证明：无论 $Y\in\operatorname{Ob}(\operatorname{Sch}/S)$、$y\in Y$ 和 $s\in\operatorname{SL}(\underline n)(Y)$ 如何选取，总能找到 $y$ 的一个“\Term{etale-neighbourhood}{平展邻域}” $Y'\to Y$，使 $s$ 在 $Y'$ 上的\Term{restriction}{限制}是 $\operatorname{GL}(\underline n)(Y')$ 中\Term{commutator}{换位子}的乘积。问题是\Term{local-scheme}{局部}的，因此可设 $Y$ 为环 $A$ 的\Term{affine}{仿射}\Term{scheme}{概形}，且 $\underline n$ 在 $Y$ 上恒等于 $n$。设 $\mathfrak p$ 是 $y$ 的\Term{prime-ideal}{素理想}，$(s_{ij})_{i,j=1,\ldots,n}$ 为 $s$ 的矩阵。对于每个 $\lambda\in A$ 和 $h,k=1,\ldots,n$、$h\ne k$，以 $c_{h,k,\lambda}$ 表示矩阵 $(c_{ij})_{i,j=1,\ldots,n}$，其中 $c_{hk}=\lambda$，且当 $(i,j)\ne(h,k)$ 时 $c_{ij}=\delta_{ij}$（$\delta_{ij}$ 是 Kronecker \Term{signature}{符号}）。对任意矩阵 $a=(a_{ij})_{i,j=1,\ldots,n}$，将 $a$ 的第 $h$ 列的 $\lambda$ 倍加到第 $k$ 列，便得到 $c_{h,k,\lambda}a$。同样，将 $a$ 的第 $k$ 行的 $\lambda$ 倍加到第 $h$ 行，便得到矩阵 $ac_{h,k,\lambda}$。特别地，
\[
 \det(c_{h,k,\lambda})=1,
\]
而且立即看出，对于每个 $\lambda,\mu\in A$，都有
\[
 c_{h,k,\lambda}c_{h,k,\mu}=c_{h,k,\lambda+\mu},\tag{6.3.1}
\]
从而得到一个\Term{representation}{表示}
\[
 c_{h,k}:\mathbf G_a\longrightarrow\operatorname{SL}(n)
 \qquad(h,k=1,\ldots,n;\ h\ne k).
\]
除非明确作相反说明，以下总假设 $c_{h,k,\lambda}$ 中的 $\lambda$（$\lambda\in A$；$h,k=1,\ldots,n$；$h\ne k$）\Term{invertible}{可逆}。证明分两步。先证明：在 $y$ 的充分小的 \Term{zariski-neighbourhood}{Zariski 邻域}中，$s$ 是矩阵 $c_{h,k,\lambda}$ 的乘积；然后证明：在 $y$ 的充分小的\Term{etale-neighbourhood}{平展邻域}中，矩阵 $c_{h,k,\lambda}$ 是 $\operatorname{GL}(n)$ 中\Term{commutator}{换位子}的乘积。

\SourcePage{567}
第一步，是在 $y$ 的一个 \Term{zariski-neighbourhood}{Zariski 邻域}中，在 $s$ 的左右两侧乘以适当的矩阵 $c_{h,k,\lambda}$，将 $s$ 化为\Term{unit}{单位}矩阵 $1=(\delta_{ij})$。先证明可用这种方法使 $s_{11}=1$。若 $s_{11}\ne1$ 且 $s_{21}\in\mathfrak p$，则要找一个矩阵 $c=c_{h,1,\lambda}$，使乘积 $s'=sc=(s'_{ij})$ 满足 $s'_{21}\notin\mathfrak p$。因为 $\det(s)=1\notin\mathfrak p$，至少存在一个指标 $i$ 使 $s_{i1}\notin\mathfrak p$。于是 $s_{21}+s_{i1}\notin\mathfrak p$，$c_{i,1,1}$ 就满足要求。因此可设 $s_{21}\notin\mathfrak p$，经过\Term{localization}{局部化}还可设 $s_{21}$ \Term{invertible}{可逆}。若 $s_{11}-1\in\mathfrak p$，则在 $s'=(s'_{ij})=sc_{1,2,1}$ 中有
\[
 s'_{11}-1=(s_{11}+s_{21})-1\notin\mathfrak p.
\]
经过\Term{localization}{局部化}，便又可设 $s_{11}-1$ \Term{invertible}{可逆}。于是 $\lambda=(1-s_{11})/s_{21}$ \Term{invertible}{可逆}，并且在 $s'=(s'_{ij})=sc_{1,2,\lambda}$ 中得到 $s'_{11}=1$。

假设 $s_{11}=1$，则可在 $y$ 的一个 \Term{zariski-neighbourhood}{Zariski 邻域}中，在 $s$ 的左侧（相应地，右侧）乘以矩阵 $c_{i,1,\lambda}$（相应地，$c_{i,1,\lambda}$），将 $s$ 化为满足
\[
 s_{11}=1,\qquad s_{1i}=s_{i1}=0\qquad(i=2,\ldots,n)
\]
的\Term{form}{形式}。然后使 $s_{22}=1$，依此类推。归纳下去，便可在 $y$ 的一个 \Term{zariski-neighbourhood}{Zariski 邻域} $U$ 中，将 $s$ 化为矩阵 $s'=(s'_{ij})$，使得当 $i<n$ 或 $j<n$ 时 $s'_{ij}=\delta_{ij}$。由构造，在 $U$ 上有 $s'=c_1sc_2$，其中 $c_1$ 和 $c_2$ 是矩阵 $c_{h,k,\lambda}$ 的乘积。因此
\[
 \det(s')=\det(s)=1,
\]
故 $s'_{nn}=1$、$s'=1$，第一步得证。

于是可设 $s$ 在 $Y$ 上已经是矩阵 $c_{h,k,\lambda}$ 的乘积。还须说明，这些矩阵在 $y$ 的一个\Term{etale-neighbourhood}{平展邻域}中是\Term{commutator}{换位子}的乘积。首先，任意两个矩阵
\SourcePage{568}
$c_{h,k,\lambda}$ 和 $c_{h,k,\mu}$ 都\Term{conjugacy}{共轭}，也就是说，存在\Term{invertible}{可逆}矩阵 $a=(a_{ij})_{i,j=1,\ldots,n}$，使得
\[
 ac_{h,k,\lambda}=c_{h,k,\mu}a.
\]
事实上，取矩阵 $(a_{ij})$ 使 $a_{hh}=1/\mu$、$a_{kk}=1/\lambda$，其余 $a_{ij}=\delta_{ij}$，就满足要求（回忆按假设 $\lambda$、$\mu$ \Term{invertible}{可逆}且 $h\ne k$）。所以，对于从 $\operatorname{GL}(n)(U)$ 到某个\Term{abelian-group}{阿贝尔群}\Term{sheaf}{层} $H$ 的任意\Term{homomorphism}{同态} $\beta$，都有
\[
 \beta(c_{h,k,\lambda})=\beta(c_{h,k,\mu}).
\]
若 $\lambda+\mu$ \Term{invertible}{可逆}，由此和 (6.3.1) 可知，对所有 $h,\beta$，都有 $\beta(c_{h,k,\lambda})=0$，所以 $c_{h,k,\lambda}$ 是\Term{commutator}{换位子}的乘积。若 $y$ 的\Term{residue-field}{剩余域} $k$ 不是\Term{prime-field}{素域} $\mathbf F_2=\{0,1\}$，则对每个\Term{invertible-element}{可逆元} $\lambda\in A$，都能找到 $y$ 的一个 \Term{zariski-neighbourhood}{Zariski 邻域} $U=\operatorname{Spec}(B)$ 和一个\Term{invertible-element}{可逆元} $\mu\in B$，使 $\lambda+\mu\in B$ 也\Term{invertible}{可逆}。因此 $c_{h,k,\lambda}$ 在 $U$ 上是\Term{commutator}{换位子}的乘积。

若 $k=\mathbf F_2$，则
\[
 V=\operatorname{Spec}\bigl(A[T]/(T^2+T+1)\bigr)
\]
是 $y$ 的一个\Term{etale-neighbourhood}{平展邻域}，其\Term{residue-field}{剩余域}不同于 $\mathbf F_2$。证明完成。

\textbf{6.5。注记。}

\textbf{6.5.1。}若 $S$ 的所有\Term{residue-field}{剩余域}都不\Term{isomorphism}{同构}于 $\mathbf F_2$，则上述证明表明，6.2.2 对 $S$ 上 Zariski \Term{sheaf}{层}构成的\Term{ringed-topos}{赋环拓扑斯}仍成立。相反，已经在 $S=\operatorname{Spec}(\mathbf F_2)$ 时，$\operatorname{SL}(2,\mathbf F_2)=\operatorname{GL}(2,\mathbf F_2)$ 就不同于 $\operatorname{GL}(2,\mathbf F_2)$ 的\Term{commutator-subgroup}{换位子群}。后者的阶为 $3$，而 $\operatorname{SL}(2,\mathbf F_2)$ 的阶为 $6$（参见 [1]）。

\textbf{6.5.2。}在任意\Term{locally-ringed-topos}{局部赋环拓扑斯} $X$ 上，\Term{generic}{一般}不能断言每个\Term{representation}{表示} $\mathbf G_m\to\mathbf G_m$ 都由一个幂给出。
\SourcePage{569}
例如，在由域 $\mathbf Q$ \Term{ringed}{赋环}的\Term{topos}{拓扑斯} $(\operatorname{Ens})$ 上，由
\[
 r\longmapsto |r|\qquad(r\in\mathbf Q^*)
\]
定义的\Term{representation}{表示} $\mathbf G_m\to\mathbf G_m$ 就不是幂映射。特别地，对任意\Term{locally-ringed-topos}{局部赋环拓扑斯} $X$ 及任意\Term{representation}{表示} $\alpha:G\to\mathbf G_m$，6.2.3 并非总成立。

\textbf{6.5.3。}不过，对任意\Term{ringed-topos}{赋环拓扑斯} $X=(C,O)$ 上与公式 (6.2.8)、(6.2.9) 和 (6.2.10) 相关联的\Term{representation}{表示}，6.2.3 仍然成立，因为这些\Term{representation}{表示}在下述意义下是“\Term{algebraic-representation}{代数的}”：

设 $e$ 是 $X$ 的\Term{final-object}{终对象}，$S=\operatorname{Spec}(O(e))$，并考虑\Term{functor}{函子}
\[
 u:C\longrightarrow(\operatorname{Sch}/S),
\]
它将每个 $Y\in\operatorname{Ob}(C)$ 对应到 $S$ 上的\Term{affine}{仿射}\Term{scheme}{概形} $\operatorname{Spec}(O(Y))$。于是，对每个正整数 $n$，$C$ 上的\Term{presheaf}{预层} $\operatorname{GL}(n)_X$ 是 $(\operatorname{Sch}/S)$ 上相应的\Term{presheaf}{预层} $\operatorname{GL}(n)_S$ 与 $u$ 的复合。称 $X$ 上的一个\Term{representation}{表示}
\[
 \alpha:\operatorname{GL}(n)\longrightarrow\mathbf G_m\tag{6.5.4}
\]
是\Term{algebraic-representation}{代数的}，若它通过 $u$ 来自 $(\operatorname{Sch}/S)$ 上群\Term{presheaf}{预层}的一个\Term{representation}{表示}
\[
 \alpha_S:\operatorname{GL}(n)_S\longrightarrow\mathbf G_{m,S},\tag{6.5.5}
\]
即对于每个 $Y\in\operatorname{Ob}(C)$ 和 $s\in\operatorname{GL}(n)(Y)=\operatorname{GL}(n)_S(u(Y))$ 都有
\[
 \alpha(s)=\alpha_S(s).
\]
当然，(6.5.4) 中的\Term{representation}{表示} $\alpha$ 完全由 $\alpha_S$ 确定。立即可知，当 $\alpha$ 是\Term{algebraic-representation}{代数的}时，(6.2.3) 在任意\Term{ringed-topos}{赋环拓扑斯}中仍成立。因此，特别地，公式 (6.2.8)、(6.2.9) 和 (6.2.10) 在任意\Term{ringed-topos}{赋环拓扑斯}上仍成立。
